% Part: computability
% Chapter: computability-theory
% Section: ce-sets

\documentclass[../../../include/open-logic-section]{subfiles}

\begin{document}

\olfileid{cmp}{thy}{ces}
\olsection{Himpunan kang Bisa Dienumerasi sacara Komputabel}

\begin{defn}
Sawijining himpunan \emph{bisa dienumerasi sacara
  komputabel} yen kosong utawa dadi citra
sawijining fungsi komputabel.
\end{defn}

\begin{history}
Himpunan kang bisa dienumerasi sacara komputabel
uga diarani \emph{bisa dienumerasi sacara
  rekursif}. Iki istilah wiwitane, lan saiki
loro-lorone lumrah digunakake, uga cekakan
``c.e.'' lan ``r.e.''
\end{history}

\begin{explain}
Coba pikirna teges definisi iki lan ngapa
istilahe cocog. Gagasane, yen $S$ minangka
citra fungsi komputabel~$f$, mula
\[
S = \{ f(0), f(1), f(2), \dots \},
\]
saengga $f$ bisa dianggep ``ngenumerasi''
unsur-unsur~$S$. Gatekna yen miturut definisi,
$f$~ora kudu fungsi kang tansah mundhak,
yaiku enumerasine ora kudu miturut urutan
mundhak. Malah $f$ uga ora kudu injektif,
yaiku ulangan ing enumerasi $f(0)$,
$f(1)$, $f(2)$, \dots{} saka~$S$ diidini.
Umpamane, fungsi konstan $f(x) = 0$
ngenumerasi himpunan $\{ 0
\}$.
\end{explain}

Saben himpunan komputabel bisa dienumerasi
sacara komputabel. Kanggo ndeleng iki,
upamane $S$~komputabel. Yen $S$ kosong,
mula miturut definisi himpunan iku bisa
dienumerasi sacara komputabel. Yen ora,
tetepna $a$ minangka unsur sembarang saka
$S$. Tetepna $f$ minangka
\[
f(x) =
\begin{cases}
x & \text{yen $\Char{S}(x) = 1$} \\
a & \text{yen ora.}
\end{cases}
\]
Mula $f$ fungsi komputabel, lan $S$
minangka citra~$f$.

\end{document}
